% !TeX encoding = UTF-8
\documentclass[variant=guia,cover=institutional,mode=print]{ua-engineering-v6-article}
% !TeX encoding = UTF-8
\UASetUniversity{Universidad Austral}
\UASetFaculty{Facultad de Ingeniería}
\UASetDepartment{Departamento de Computación}
\UASetCampus{Pilar}
\UASetCareer{Ingeniería Informática}
\UASetCommission{Comisión A}
\UASetProfessor{Equipo docente}
\UASetSemester{Segundo cuatrimestre 2026}
\UASetCourse{Sistemas Distribuidos}
\UASetDate{2026}

\UASetSemester{Segundo cuatrimestre 2026}
\UASetDate{2026}
\usetikzlibrary{positioning,arrows.meta,calc}

\UASetClassNumber{06}
\UASetClassTitle{Storage distribuido y almacenamiento permanente}
\UASetDocKind{Ficha de trabajo en aula}
\title{Clase 06 --- Ficha de aula: persistencia, concurrencia y rendimiento bajo fallas}
\author{\UAProfessor}
\date{\UADate}

\begin{document}
\setlength{\emergencystretch}{2em}
\maketitle

\section{Contexto y vocabulario inicial}

\textbf{Pregunta central:} ¿cómo convierte un sistema distribuido una operación en estado permanente, recuperable y consultable?

Complete durante la apertura:

\begin{center}
\begin{tabular}{p{0.25\textwidth}p{0.60\textwidth}}
\toprule
Concepto & Definición en sus palabras \\
\midrule
Capa distribuida & \rule{0.58\textwidth}{0.4pt} \\[0.45cm]
Motor de almacenamiento & \rule{0.58\textwidth}{0.4pt} \\[0.45cm]
WAL & \rule{0.58\textwidth}{0.4pt} \\[0.45cm]
MVCC & \rule{0.58\textwidth}{0.4pt} \\[0.45cm]
LSM-tree & \rule{0.58\textwidth}{0.4pt} \\[0.45cm]
COMMIT & \rule{0.58\textwidth}{0.4pt} \\[0.45cm]
X durable & \rule{0.58\textwidth}{0.4pt} \\[0.45cm]
\bottomrule
\end{tabular}
\end{center}

\begin{uainfo}
\textbf{Caso CampusShop.} ID lógico: 9F2A. Filas: \texttt{purchases} e \texttt{inventory}. Página didáctica: P17. Entrada de índice: I5.
\end{uainfo}

\newpage
\section{Timeline WAL}

Ubique: request 9F2A, preparación, \texttt{WAL UPDATE}, páginas sucias, \texttt{COMMIT}, flush hasta \texttt{commitLSN}, ACK, crash y REDO.

\begin{center}
\begin{tikzpicture}[x=1.0cm,y=1.0cm,>=Latex,every node/.style={font=\small}]
\foreach \y/\lab in {4/Cliente,3/Servicio,2/Memoria,1/WAL,0/Páginas}{
  \node[anchor=east] at (0.75,\y){\textbf{\lab}};
  \draw[UAIngRule,line width=0.9pt] (0.95,\y)--(10.2,\y);
}
\draw[->,very thick,UAIngNavy] (0.95,-0.7)--(10.4,-0.7) node[right]{tiempo};
\foreach \x in {1.5,2.5,3.5,4.5,5.5,6.5,7.5,8.5,9.5}{
  \draw[UAIngMuted] (\x,-0.58)--(\x,-0.82);
}
\end{tikzpicture}
\end{center}

\noindent\textbf{Garantía después de recovery:}\\[0.18cm]\rule{\textwidth}{0.4pt}

\noindent\textbf{Punto más temprano de ACK legítimo:}\\[0.18cm]\rule{\textwidth}{0.4pt}

\noindent\textbf{¿El flush durable comienza con WAL UPDATE?}\\[0.18cm]\rule{\textwidth}{0.4pt}

\noindent\textbf{Prefijo que debe quedar durable:}\\[0.18cm]\rule{\textwidth}{0.4pt}

\newpage
\section{Crashes y recuperación}

\begin{center}
\begin{tabular}{p{0.31\textwidth}p{0.22\textwidth}p{0.31\textwidth}}
\toprule
Punto de crash & ¿Qué sabe el cliente? & Estado después de recovery \\
\midrule
Antes de COMMIT & & \\[1.2cm]
COMMIT agregado, no flushed & & \\[1.2cm]
Flush durable, antes de ACK & & \\[1.2cm]
ACK, antes de P17 & & \\[1.2cm]
\bottomrule
\end{tabular}
\end{center}

\textbf{Papel de 9F2A ante un retry:}

\vspace{1.0cm}

\textbf{Si las páginas conservan el estado anterior, complete:}

\[
\text{estado recuperado} = \rule{3.1cm}{0.4pt} + \operatorname{REDO}(\rule{3.1cm}{0.4pt})
\]

\textbf{Registros WAL de modificación contienen:} \rule{0.48\textwidth}{0.4pt}

\textbf{El registro COMMIT indica:} \rule{0.53\textwidth}{0.4pt}

\vspace{0.4cm}

\textbf{Diferencia entre checkpoint, réplica y backup:}

\vspace{1.7cm}

\newpage
\section{Timeline MVCC}

Ubique: $v_1$, inicio de $T_R$, snapshot $S_R$, primera lectura, creación de $v_2$, COMMIT de $T_W$, segunda lectura, lector nuevo, fin de $T_R$ y reclamación de $v_1$.

\begin{center}
\begin{tikzpicture}[x=1.0cm,y=1.05cm,>=Latex,every node/.style={font=\small}]
\foreach \y/\lab in {3/{Lector $T_R$},2/{Escritor $T_W$},1/Versiones,0/{Lectura visible}}{
  \node[anchor=east] at (0.75,\y){\textbf{\lab}};
  \draw[UAIngRule,line width=0.9pt] (0.95,\y)--(10.2,\y);
}
\draw[->,very thick,UAIngNavy] (0.95,-0.7)--(10.4,-0.7) node[right]{tiempo};
\end{tikzpicture}
\end{center}

\textbf{Snapshot no significa:} \rule{0.58\textwidth}{0.4pt}

\textbf{Regla de visibilidad en una frase:}

\vspace{1.0cm}

\textbf{Costo de una transacción larga:}

\vspace{1.0cm}

\textbf{Por qué MVCC no implica serializabilidad:}

\vspace{1.0cm}

\newpage
\section{Transición: MVCC, árbol B y ruta física}

Complete las dos preguntas:

\begin{center}
\begin{tabular}{p{0.29\textwidth}p{0.58\textwidth}}
\toprule
Mecanismo & Pregunta que responde \\
\midrule
MVCC & \rule{0.55\textwidth}{0.4pt} \\[0.6cm]
Árbol B / índice & \rule{0.55\textwidth}{0.4pt} \\[0.6cm]
\bottomrule
\end{tabular}
\end{center}

\textbf{Regla de transición:} el índice localiza \rule{0.35\textwidth}{0.4pt}; MVCC decide \rule{0.30\textwidth}{0.4pt}.

\vspace{0.35cm}

Dibuje una ruta desde raíz hasta hoja y marque dónde se comprueba visibilidad:

\begin{center}
\begin{tikzpicture}[x=1cm,y=1cm,>=Latex,every node/.style={font=\small}]
\node[draw=UAIngNavy,rounded corners,minimum width=2.0cm,minimum height=0.7cm] (r) at (5,3.0) {raíz};
\node[draw=UAIngNavy,rounded corners,minimum width=2.0cm,minimum height=0.7cm] (i1) at (2.3,1.7) {página interna};
\node[draw=UAIngNavy,rounded corners,minimum width=2.0cm,minimum height=0.7cm] (i2) at (7.7,1.7) {página interna};
\node[draw=UAIngOrange,rounded corners,minimum width=2.0cm,minimum height=0.7cm] (l1) at (1.2,0.2) {hoja};
\node[draw=UAIngOrange,rounded corners,minimum width=2.0cm,minimum height=0.7cm] (l2) at (3.5,0.2) {hoja};
\node[draw=UAIngOrange,rounded corners,minimum width=2.0cm,minimum height=0.7cm] (l3) at (6.5,0.2) {hoja};
\node[draw=UAIngOrange,rounded corners,minimum width=2.0cm,minimum height=0.7cm] (l4) at (8.8,0.2) {hoja};
\draw[->,UAIngNavy,thick] (r)--(i1);
\draw[->,UAIngNavy,thick] (r)--(i2);
\draw[->,UAIngNavy,thick] (i1)--(l1);
\draw[->,UAIngNavy,thick] (i1)--(l2);
\draw[->,UAIngNavy,thick] (i2)--(l3);
\draw[->,UAIngNavy,thick] (i2)--(l4);
\end{tikzpicture}
\end{center}

\textbf{Fan-out significa:} \rule{0.60\textwidth}{0.4pt}

\textbf{Por qué poca altura ayuda:} \rule{0.56\textwidth}{0.4pt}

\newpage
\section{EXPLAIN: comparar evidencia}

Consulta:

\begin{verbatim}
SELECT id, total, confirmed_at
FROM purchases
WHERE student_id = 8421
  AND status = 'CONFIRMED'
ORDER BY confirmed_at DESC
LIMIT 20;
\end{verbatim}

\begin{center}
\begin{tabular}{p{0.28\textwidth}p{0.25\textwidth}p{0.25\textwidth}}
\toprule
Evidencia & Sin índice & Con índice compuesto \\
\midrule
Acceso inicial & & \\[0.7cm]
Filas descartadas & & \\[0.7cm]
Buffers & & \\[0.7cm]
Sort separado & & \\[0.7cm]
Trabajo evitado & & \\[0.9cm]
Costo agregado a writes & & \\[0.9cm]
\bottomrule
\end{tabular}
\end{center}

\textbf{DDL del índice:}

\vspace{1.1cm}

\textbf{Razón válida para preferir Seq Scan:}

\vspace{1.0cm}

\newpage
\section{LSM-tree, merge y filtro de Bloom}

Complete el camino de escritura:

\[
(k,v) \rightarrow \rule{1.3cm}{0.4pt} \rightarrow
\rule{1.5cm}{0.4pt} \rightarrow \rule{2.0cm}{0.4pt}
\rightarrow \rule{1.2cm}{0.4pt} \rightarrow \text{SSTable L0}.
\]

\begin{center}
\begin{tabular}{p{0.25\textwidth}p{0.58\textwidth}}
\toprule
Término & Explicación \\
\midrule
SSTable & \\[0.8cm]
Merge ordenado & \\[0.8cm]
Tombstone & \\[0.8cm]
Compactación & \\[0.8cm]
\bottomrule
\end{tabular}
\end{center}

\textbf{Filtro de Bloom:}

\begin{itemize}
\item ``definitivamente ausente'' significa: \rule{0.54\textwidth}{0.4pt}
\item ``posiblemente presente'' significa: \rule{0.54\textwidth}{0.4pt}
\item falso positivo: \rule{0.35\textwidth}{0.4pt}
\item falso negativo admisible: Sí / No. Justificación: \rule{0.33\textwidth}{0.4pt}
\end{itemize}

Reemplace ``archivo imposible'' por una formulación precisa:

\vspace{0.8cm}

\newpage
\section{Compactación y write stall}

Datos:

\begin{itemize}
\item ingreso: 90 MB/s;
\item capacidad de compactación: 60 MB/s;
\item backlog: 18 GB;
\item slowdown: 27 GB;
\item stop: 36 GB;
\item L0: 28/36;
\item p99: creciente.
\end{itemize}

\textbf{Déficit:} \rule{0.28\textwidth}{0.4pt}

\textbf{Tiempo ideal hasta slowdown:} \rule{0.28\textwidth}{0.4pt}

\textbf{Tiempo ideal hasta stop:} \rule{0.28\textwidth}{0.4pt}

\begin{center}
\begin{tabular}{p{0.35\textwidth}p{0.43\textwidth}}
\toprule
Hipótesis & Métrica para verificarla \\
\midrule
1. & \\[0.9cm]
2. & \\[0.9cm]
3. & \\[0.9cm]
4. & \\[0.9cm]
\bottomrule
\end{tabular}
\end{center}

\textbf{Un stall protege cuando:}

\vspace{0.8cm}

\textbf{Un stall revela incapacidad cuando:}

\vspace{0.8cm}

\newpage
\section{Decisión final y exit ticket}

\begin{center}
\begin{tabular}{p{0.24\textwidth}p{0.31\textwidth}p{0.31\textwidth}}
\toprule
Dimensión & CampusShop & Telemetría del campus \\
\midrule
Patrón de carga & & \\[0.7cm]
Falla crítica & & \\[0.7cm]
Garantía & & \\[0.7cm]
Estructura candidata & & \\[0.7cm]
Costo desplazado & & \\[0.7cm]
Métrica principal & & \\[0.7cm]
Experimento & & \\[0.9cm]
\bottomrule
\end{tabular}
\end{center}

\subsection*{Exit ticket}

\begin{enumerate}
\item Diferencia entre aceptado, escrito y durable.
\item Costo oculto de MVCC y una métrica para observarlo.
\item Diferencia entre \texttt{EXPLAIN} y \texttt{EXPLAIN ANALYZE}.
\item Señales de que una LSM acumula trabajo pendiente.
\end{enumerate}

\end{document}
